\documentclass[11pt]{article}
\usepackage[margin=1in]{geometry}
\usepackage{amsmath,amssymb,amsthm}
\usepackage{array}
\usepackage{parskip}
\usepackage{booktabs}

\newtheorem{theorem}{Theorem}
\newtheorem{lemma}[theorem]{Lemma}
\newtheorem{corollary}[theorem]{Corollary}
\newtheorem{proposition}[theorem]{Proposition}
\theoremstyle{definition}
\newtheorem{definition}[theorem]{Definition}
\theoremstyle{remark}
\newtheorem{remark}[theorem]{Remark}

\newcommand{\N}{\mathbb{N}}
\newcommand{\Ap}{\operatorname{Ap}}
\newcommand{\gen}[1]{\langle #1 \rangle}

\title{Disproof of conjecture \texttt{00000000588}}
\author{earthking11}
\date{2026-09-14}

\begin{document}
\maketitle

\section*{Verdict: FALSE}

We disprove conjecture \texttt{00000000588} as stated. The refutation is
unconditional, completely elementary, and consists of a single explicit
counterexample together with the correct closed form:
\[
t\bigl(\gen{n, n+1, \dots, 2n-1}\bigr) = n-1 \qquad (n \ge 2),
\]
which equals the conjectured value $\lceil n/2 \rceil$ only at $n = 2$ and
$n = 3$. Already at $n = 4$, the smallest case not covered by the coincidence,
we have
\[
S = \gen{4,5,6,7}, \qquad \Ap(S,4) = \{0,5,6,7\}, \qquad t(S) = 3
\;\neq\; 2 = \lceil 4/2 \rceil .
\]
The value $t = 3$ is not merely computed: it is formalised in Lean~4 with a
self-contained, import-free, axiom-free proof (Section~\ref{sec:lean}).

\section{The conjecture, quoted verbatim}

From \texttt{conjectures/00000000588.md}:

\begin{quote}
\textbf{English.} Definition: The type $t(S)$ (the number of maximal Apéry
elements). Conjecture: $t(\gen{n, n+1, \dots, 2n-1}) = \lceil n/2 \rceil$
(a closed form for the type of interval semigroups).
\end{quote}

The conjecture is filed in both English and Chinese; the Chinese original is
reproduced verbatim (in Unicode) in \texttt{README.md}, and the two versions are
identical in content. The PDF font does not embed CJK glyphs, so the Chinese is
kept in the markdown copy rather than duplicated here.

\section{Definitions and conventions}

We pin every notion used by the conjecture.

\begin{definition}[Numerical semigroup]\label{def:semigroup}
A \emph{numerical semigroup} is a submonoid $S \subseteq \N$ under addition with
finite complement $\N \setminus S$. For $A \subseteq \N$ we write $\gen{A}$ for
the submonoid generated by $A$, i.e.\ the set of finite sums of elements of $A$
(including the empty sum $0$).
\end{definition}

\begin{definition}[Interval semigroup]
For $n \ge 1$, the \emph{interval semigroup} is
$S_n = \gen{n, n+1, \dots, 2n-1}$.
\end{definition}

\begin{definition}[Multiplicity]
The \emph{multiplicity} of $S \neq \{0\}$ is its smallest positive element,
$m(S) = \min(S \setminus \{0\})$.
\end{definition}

\begin{definition}[Apéry set]\label{def:apery}
For a numerical semigroup $S$ and $m \in S \setminus \{0\}$,
\[
\Ap(S,m) = \{\, s \in S \;:\; s - m \notin S \,\},
\]
where $s - m$ is computed in $\mathbb{Z}$ (so $0 \in \Ap(S,m)$ because
$-m \notin S$).
\end{definition}

The conjecture says ``the Apéry elements'' without naming a modulus. The
standard convention, which we adopt throughout, is the \emph{multiplicity}
$m = m(S)$, the smallest positive element; for $S_n$ that is $m = n$. Section
\ref{sec:robust} records that the verdict does not depend on this choice.

\begin{definition}[Apéry order]\label{def:order}
On $S$ put $x \le_S y$ iff $y - x \in S$. This is a partial order on $S$. Its
strict part is $x <_S y$ iff $x < y$ and $y - x \in S$. A \emph{maximal Apéry
element} is an element of $\Ap(S,m)$ that is maximal for $\le_S$ restricted to
$\Ap(S,m)$.
\end{definition}

\begin{definition}[Type]\label{def:type}
The \emph{type} $t(S)$ is the number of maximal Apéry elements of $S$.
\end{definition}

\begin{remark}
The maximal Apéry elements are in bijection with the pseudo-Frobenius numbers
of $S$; equivalently, $t(S) = |\mathrm{PF}(S)|$ where
$\mathrm{PF}(S) = \{\, x \notin S : x + s \in S \ \forall s \in S\setminus\{0\}
\,\}$. We use this only as a cross-check (Section~\ref{sec:robust}).
\end{remark}

\section{The counterexample $n = 4$}\label{sec:n4}

Throughout this section $S = \gen{4,5,6,7}$ and we use the multiplicity $4$.

\begin{lemma}\label{lem:S}
$S = \{0\} \cup [4, \infty) = \{0\} \cup \{\, m \in \N : m \ge 4 \,\}$.
\end{lemma}

\begin{proof}
The inclusion $\subseteq$ is clear since every generator is $\ge 4$ and $0$
is the empty sum. For the converse, use strong induction on $m \ge 4$. If
$4 \le m \le 7$, then $m$ is one of the generators $4,5,6,7$, hence $m \in S$.
If $m \ge 8$, then $4 \in S$ and $m - 4 \ge 4$; by induction $m - 4 \in S$, so
$m = 4 + (m-4) \in S$ as $S$ is closed under addition.
\end{proof}

\begin{proposition}\label{prop:apery4}
$\Ap(S,4) = \{0, 5, 6, 7\}$.
\end{proposition}

\begin{proof}
By Lemma~\ref{lem:S}, $S = \{0\} \cup [4,\infty)$. Now
\[
\Ap(S,4) = \{\, s \in S : s - 4 \notin S \,\}.
\]
We test the elements of $S$:
\begin{itemize}
\item $s = 0$: $0 - 4 = -4 \notin S$, so $0 \in \Ap(S,4)$.
\item $s = 4$: $4 - 4 = 0 \in S$, so $4 \notin \Ap(S,4)$.
\item $s \in \{5,6,7\}$: $s - 4 \in \{1,2,3\}$, and $1,2,3 \notin S$ because
  $S$ contains $0$ and all integers $\ge 4$; so $5,6,7 \in \Ap(S,4)$.
\item $s \ge 8$: then $s - 4 \ge 4$, so $s - 4 \in S$ and
  $s \notin \Ap(S,4)$.
\end{itemize}
Hence $\Ap(S,4) = \{0,5,6,7\}$.
\end{proof}

\begin{proposition}\label{prop:t4}
$t(S) = 3$, whereas $\lceil 4/2 \rceil = 2$.
\end{proposition}

\begin{proof}
We determined $\Ap(S,4) = \{0,5,6,7\}$ in Proposition~\ref{prop:apery4}. On
this set the order $\le_S$ has strict part $x <_S y$ iff $x < y$ and
$y - x \in S$. We check maximality.

\emph{The three nonzero elements are pairwise incomparable.} For distinct
$x, y \in \{5,6,7\}$ with $x < y$ we have $y - x \in \{1,2\}$ (indeed
$6-5 = 1$, $7-6 = 1$, $7-5 = 2$), and $1, 2 \notin S$. Hence
$\neg(x <_S y)$; as $x < y$ was arbitrary, no nonzero Apéry element lies below
another.

\emph{Each of $5, 6, 7$ is maximal.} By the previous paragraph there is no
element of $\Ap(S,4)$ strictly above it, so each is maximal.

\emph{$0$ is not maximal.} $0 <_S 5$ because $0 < 5$ and $5 - 0 = 5 \in S$.

Therefore the maximal Apéry elements are exactly $\{5,6,7\}$ and
\[
t(S) = |\{5,6,7\}| = 3 \neq 2 = \lceil 4/2 \rceil . \qedhere
\]
\end{proof}

\begin{corollary}
Conjecture \texttt{00000000588} is false.
\end{corollary}

\begin{proof}
The conjecture asserts $t(S_n) = \lceil n/2 \rceil$ for the interval
semigroups. For $n = 4$ Proposition~\ref{prop:t4} gives $t(S_4) = 3$ while
$\lceil 4/2 \rceil = 2$; the universally quantified closed form fails at
$n = 4$.
\end{proof}

\section{The true closed form}\label{sec:closedform}

The counterexample is not an isolated anomaly; the correct closed form is
$n - 1$.

\begin{theorem}\label{thm:main}
For every $n \ge 2$, $S_n = \gen{n, n+1, \dots, 2n-1}$ satisfies
\[
S_n = \{0\} \cup [n,\infty), \qquad
\Ap(S_n, n) = \{0\} \cup \{n+1, n+2, \dots, 2n-1\},
\]
and every element of $\{n+1, \dots, 2n-1\}$ is maximal, while $0$ is not.
Consequently
\[
t(S_n) = n - 1 .
\]
\end{theorem}

\begin{proof}
\emph{Step 1: $S_n = \{0\} \cup [n,\infty)$.} Every generator lies in
$[n, 2n-1] \subseteq [n,\infty)$ and $0$ is the empty sum, giving
$\subseteq$. Conversely, let $m \ge n$. If $m \le 2n-1$ then $m$ is a
generator. If $m \ge 2n$, then $m - n \ge n$ and $m - n \le m$; by strong
induction on $m$, $m - n \in S_n$, whence $m = n + (m-n) \in S_n$.

\emph{Step 2: the Apéry set.} Let $s \in S_n$. If $s = 0$, then
$s - n = -n \notin S_n$, so $s \in \Ap(S_n,n)$. If $n \le s \le 2n-1$ and
$s = n$, then $s - n = 0 \in S_n$, so $s \notin \Ap(S_n,n)$; if
$n+1 \le s \le 2n-1$, then $s - n \in [1, n-1]$, and $[1,n-1]$ contains no
element of $S_n = \{0\} \cup [n,\infty)$ (because $n \ge 2$, so
$1 \le s-n \le n-1 < n$), hence $s \in \Ap(S_n,n)$. Finally, if $s \ge 2n$,
then $s - n \ge n$, so $s - n \in S_n$ by Step 1 and
$s \notin \Ap(S_n,n)$. Therefore
$\Ap(S_n,n) = \{0\} \cup \{n+1,\dots,2n-1\}$.

\emph{Step 3: maximality.} Let $x, y \in \Ap(S_n,n)$ with $x < y$. Then
$y - x \le (2n-1) - (n+1) = n - 2$ and $y - x \ge 1$. Since $S_n$ contains no
integer in $[1, n-1]$ (Step 1) and $n \ge 2$ gives $n - 2 \le n - 1$, we get
$y - x \notin S_n$, so $\neg(x <_S y)$. Hence the $n-1$ elements
$n+1, \dots, 2n-1$ are pairwise incomparable, and each is maximal. Moreover
$0 <_S n+1$ because $(n+1) - 0 = n+1 \in S_n$, so $0$ is not maximal.
Therefore the maximal Apéry elements are exactly $n+1, \dots, 2n-1$ and
\[
t(S_n) = (2n-1) - (n+1) + 1 = n - 1 . \qedhere
\]
\end{proof}

\begin{corollary}[Agreement with the conjecture]\label{cor:agree}
$t(S_n) = n - 1$ equals $\lceil n/2 \rceil$ if and only if $n \in \{2,3\}$.
\end{corollary}

\begin{proof}
For $n = 2$: $n - 1 = 1 = \lceil 2/2 \rceil$. For $n = 3$:
$n - 1 = 2 = \lceil 3/2 \rceil$. For $n \ge 4$ we have
$n - 1 > \lceil n/2 \rceil$: indeed $n - 1 - \lceil n/2 \rceil
= \lfloor (n-2)/2 \rfloor \ge 1$. So the two functions differ at every
$n \ge 4$, and the conjecture fails for all of them.
\end{proof}

\section{The table}\label{sec:table}

The following values are produced by \texttt{reproduce.py} from an exact
membership test (Section~\ref{sec:repro}); the computation for $n = 4$ is the
one verified in Lean.

\begin{center}
\begin{tabular}{r l l r r c}
\toprule
$n$ & $S_n$ & $\Ap(S_n, n)$ & $t(S_n)$ & $\lceil n/2 \rceil$ & agrees? \\
\midrule
$2$ & $\gen{2,3}$       & $\{0,3\}$                   & $1$ & $1$ & yes \\
$3$ & $\gen{3,4,5}$     & $\{0,4,5\}$                 & $2$ & $2$ & yes \\
$4$ & $\gen{4,5,6,7}$   & $\{0,5,6,7\}$               & $3$ & $2$ & \textbf{no} \\
$5$ & $\gen{5,\dots,9}$ & $\{0,6,7,8,9\}$             & $4$ & $3$ & \textbf{no} \\
$6$ & $\gen{6,\dots,11}$& $\{0,7,\dots,11\}$          & $5$ & $3$ & \textbf{no} \\
$7$ & $\gen{7,\dots,13}$& $\{0,8,\dots,13\}$          & $6$ & $4$ & \textbf{no} \\
$8$ & $\gen{8,\dots,15}$& $\{0,9,\dots,15\}$          & $7$ & $4$ & \textbf{no} \\
$9$ & $\gen{9,\dots,17}$& $\{0,10,\dots,17\}$         & $8$ & $5$ & \textbf{no} \\
\bottomrule
\end{tabular}
\end{center}

The table is the closed form $t = n - 1$ against the conjectured
$\lceil n/2 \rceil$: the two agree only at $n = 2,3$ and separate for every
$n \ge 4$. The script additionally verifies $n = 10, 11, 12$.

\section{Robustness}\label{sec:robust}

We tried to save the conjecture by varying each convention; none succeeds.

\paragraph{Choice of Apéry modulus.} The type is independent of the nonzero
modulus used. For $n = 4$ and each $m \in \{4,5,6,7,9,13\}$ we computed
$\Ap(S,m) = \{s \in S : s - m \notin S\}$ and its maximal elements:

\begin{center}
\begin{tabular}{r l l r}
\toprule
modulus $m$ & $\Ap(S,m)$ & maximal elements & count \\
\midrule
$4$  & $\{0,5,6,7\}$                   & $\{5,6,7\}$   & $3$ \\
$5$  & $\{0,4,6,7,8\}$                 & $\{6,7,8\}$   & $3$ \\
$6$  & $\{0,4,5,7,8,9\}$               & $\{7,8,9\}$   & $3$ \\
$7$  & $\{0,4,5,6,8,9,10\}$            & $\{8,9,10\}$  & $3$ \\
$9$  & $\{0,4,5,6,7,8,10,11,12\}$      & $\{10,11,12\}$& $3$ \\
$13$ & $\{0,4,\dots,12,14,15,16\}$     & $\{14,15,16\}$& $3$ \\
\bottomrule
\end{tabular}
\end{center}

Every modulus gives $3$ maximal elements, and the maximal set is always the
translate by $m$ of the pseudo-Frobenius set
$\mathrm{PF}(S) = \{1,2,3\}$: indeed $\{5,6,7\} - 4 = \{1,2,3\}$,
$\{6,7,8\} - 5 = \{1,2,3\}$, and so on. Thus the type is $3$ and the
pseudo-Frobenius numbers are $\{1,2,3\}$, of size $3$ as well.

\paragraph{Rounding conventions.} For $n = 4$ all natural readings of
``$n/2$'' give $2$: $\lceil 4/2 \rceil = 2$, $\lfloor 4/2 \rfloor = 2$,
$\operatorname{round}(4/2) = 2$. Since $t = 3$, none of them restores the
claim. More generally, replacing $\lceil \cdot \rceil$ by $\lfloor \cdot
\rfloor$ or by rounding leaves $n = 5$ ($t = 4$, all roundings of $5/2$ are
$2$ or $3$) as a counterexample too.

\paragraph{Other readings.} The conjecture is a single closed-form identity
for a universally quantified $n$; it is not an asymptotic or an inequality, so
there is no constant or error term to adjust. Either the identity
$t(S_n) = \lceil n/2 \rceil$ holds for all $n$ (it does not: $n = 4$) or the
conjecture is false.

\section{Formalisation}\label{sec:lean}

The value $t(S_4) = 3$ is formalised in the Lean~4 project under
\texttt{lean4/}. \texttt{Main.lean} has \textbf{no imports at all} (not even
\texttt{Std}), uses no Mathlib and no \texttt{sorry}; it defines
\[
\texttt{inS}\ x := \exists a\,b\,c\,d,\ x = 4a+5b+6c+7d, \qquad
\texttt{inAp}\ x := x = 0 \lor x = 5 \lor x = 6 \lor x = 7,
\]
\[
\texttt{ltS}\ x\,y := x < y \land \texttt{inS}\,(y-x), \qquad
\texttt{isMax}\ x := \texttt{inAp}\ x \land \forall y,\ \texttt{inAp}\ y \to \neg\,\texttt{ltS}\ x\,y,
\]
and proves
\begin{quote}\ttfamily\small
theorem not\_inS\_1 : \(\neg\) inS 1 \quad (and likewise 2, 3)\\
theorem not\_isMax\_0 : \(\neg\) isMax 0\\
theorem isMax\_5 : isMax 5 \quad (and likewise 6, 7)\\
theorem maximals\_correct :\\
\hspace*{1em}\((\forall x,\ x \in maximals \to isMax\ x)\)\\
\hspace*{1em}\(\wedge\ (\forall x,\ isMax\ x \to x \in maximals)\)\\
theorem counterexample\_n4 : t = 3 \(\wedge\) t \(\ne\) 2
\end{quote}
Here \texttt{maximals} is the list \texttt{[5,6,7]} and
\texttt{t := maximals.length}, so \texttt{t = 3} holds by \texttt{rfl} and
\texttt{t \(\ne\) 2} by \texttt{decide}. The file also formalises $n = 3$
(\texttt{conjecture\_holds\_n3 : t3 = 2}) and $n = 5$
(\texttt{counterexample\_n5 : t5 = 4 \(\wedge\) t5 \(\ne\) 3}).
\texttt{Check.lean} runs \texttt{\#print axioms} on every theorem: the main
theorems depend on \emph{no axioms}, the rest only on the core axioms
\texttt{propext} and \texttt{Quot.sound}, and none on \texttt{sorryAx}. See
\texttt{lean4/README.md}.

\section{Reproducibility}\label{sec:repro}

\texttt{reproduce.py} uses the Python~3 standard library only. For each
$n = 2, \dots, 12$ it computes membership in $S_n$ by exact dynamic
programming (unbounded coin change: $x$ is representable iff $x = 0$ or
$x - g$ is representable for some generator $g \le x$), forms
$\Ap(S_n,n) = \{s \in S_n : s - n \notin S_n\}$, finds the maximal elements
under the Apéry order, and reports $t(S_n)$ and $\lceil n/2 \rceil$. It asserts
$t(S_4) = 3 \ne 2$, that $t(S_n) = n-1$ for every $n \ge 2$, and the robustness
checks of Section~\ref{sec:robust}, printing \texttt{PASS} and exiting $0$ iff
every check holds.

\begin{quote}\ttfamily\small
\$ python3 reproduce.py\\
... \textit{(Apéry sets, maximal elements, the $n = 2..12$ table,
robustness, all checks)}\\
PASS: all checks verified; conjecture 00000000588 is FALSE.
\end{quote}

\section*{Status against the submission rules}

Rule 3 requires a LaTeX source, a PDF document, and a Lean~4 project.

\begin{itemize}\sloppy
\item \textbf{LaTeX source} --- \texttt{main.tex}, a standalone \texttt{article}
  using \texttt{geometry}, \texttt{amsmath}, \texttt{amssymb}, \texttt{amsthm},
  \texttt{array}, \texttt{parskip}, \texttt{booktabs}; compiles with
  \texttt{tectonic main.tex}.
\item \textbf{PDF} --- at \texttt{build/main.pdf}, produced by
  \texttt{tectonic --outdir build main.tex}.
\item \textbf{Lean~4 project} --- \texttt{lean4/} with \texttt{lean-toolchain}
  pinned to \texttt{leanprover/lean4:v4.33.1}, \texttt{lakefile.toml}
  (library \texttt{Main}, no dependencies), \texttt{Main.lean} (no imports),
  \texttt{Check.lean} (\texttt{\#print axioms}), and \texttt{lean4/README.md}.
  \texttt{lake build} exits $0$; the audit reports no \texttt{sorryAx}.
\end{itemize}

\end{document}
